\section{Related Work}\label{sec:soa}

Our research bridges two bodies of literature: high-level GPU programming models and algorithmic optimization of automata processing.

\subsection{High-Level GPU Abstractions}

The complexity of CUDA has driven research into Domain-Specific Languages (DSLs).
\textbf{Triton}~\cite{triton} simplifies writing high-performance neural network kernels through a tile-based semantics where the compiler manages memory coalescing and shared memory allocation.
While Tillet et al.\ demonstrated Triton's effectiveness for matrix multiplications, their analysis was confined to dense linear algebra. Our work extends this evaluation to irregular, latency-bound algorithms.

\subsection{Parallel FSM Processing on GPUs}

Mapping Finite State Machines to GPUs is difficult due to sequential data dependencies.

\subsubsection{Speculative Approaches}
Mytkowicz et al.~\cite{mytkowicz2014} proposed enumeration-based parallelism, speculatively executing the FSM for all possible starting states. While effective on CPUs with SIMD, this approach has high overhead for single-stream processing.

\subsubsection{Memory-Centric Approaches}
Yaneva et al.~\cite{yaneva2017} analyzed memory layouts for FSM traversal, contrasting CSR against dense Transition Matrices. Our Triton implementation builds on these findings with both Bitmap and CSR variants.

\subsubsection{Worklist-Based Approaches}
Ge et al.\ proposed \textbf{ngAP}~\cite{ngap}, a non-blocking queue-based architecture using warp-level primitives for load balancing. We include ngAP-inspired implementations as a baseline, showing that the required hardware primitives (\texttt{warp\_shuffle}, fine-grained shared memory) are inexpressible in Triton.

\subsubsection{Bit-Parallel Approaches}
Recent work has explored bit-parallel methods that transform regex into bitstream programs.
\textbf{HybridSA}~\cite{hybridsa} (OOPSLA 2024) demonstrated CPU-GPU parallel engines using bit parallelism for NFA simulation, reducing irregular memory accesses.
\textbf{BitGen}~\cite{bitgen} (MICRO 2025) introduced interleaved bitstream execution, keeping intermediate data in registers rather than global memory. BitGen reports 19.5$\times$ speedup over ngAP by transforming the workload from memory-bound to compute-bound.
We evaluate BitGen in our benchmark suite to understand how bit-parallel approaches compare to traditional graph-traversal methods across different automata structures.
